\chapter{Complete Blood Count (CBC)}

\section*{What Is This Test?}
A complete blood count (CBC) is a group of tests that check the cells that circulate in your blood. It is one of the most commonly ordered blood tests. It gives information about the types and numbers of cells in the blood, including red blood cells (which carry oxygen), white blood cells (which fight infection), and platelets (which help blood clot). The CBC measures: hemoglobin (the protein in red blood cells that carries oxygen), hematocrit (the percentage of blood volume occupied by red blood cells), red blood cell count, white blood cell count with differential (the percentage and absolute number of each type of white blood cell: neutrophils, lymphocytes, monocytes, eosinophils, and basophils), platelet count, and red blood cell indices (MCV, MCH, MCHC, RDW). The CBC is used to diagnose and monitor many conditions, including anemia, infection, leukemia, lymphoma, clotting disorders, and bone marrow disorders \cite{wallach2022interpretation}.

\section*{Alternative Names}
\begin{itemize}
  \item CBC
  \item Complete Blood Count
  \item Blood Count
  \item Hemogram
  \item CBC with Differential
\end{itemize}
\section*{Principle}
The CBC is done on an automated hematology analyzer. It uses a combination of impedance counting, light scatter, conductivity, and cytochemistry. Red blood cells and platelets are counted and sized with electrical impedance. Cells pass through a small opening, and each cell produces an electrical pulse that matches its size. White blood cells are sorted using reagents that lyse red blood cells and stain white blood cells, followed by flow cytometry with laser light scatter analysis. Different manufacturers use different technologies. Beckman Coulter uses VCS (volume, conductivity, scatter) technology. Siemens uses fluorescent flow cytometry with proprietary reagents. Sysmex uses fluorescence flow cytometry with semiconductor lasers. Hemoglobin is measured by spectrophotometry after red blood cell lysis. Hematocrit is calculated from the red blood cell count and MCV. Platelet count is measured by impedance or optical methods. Newer analyzers can also report the immature platelet fraction (IPF).

\section*{History}
The CBC has its roots in the late 19th century. Paul Ehrlich developed staining techniques for blood cells in the 1870s-1880s. Manual counting with a hemocytometer (designed by Karl Thoma in 1878) was the standard method for decades. Wallace and Joseph Coulter patented the Coulter principle (electrical impedance cell counting) in 1953. It revolutionized hematology. The first automated hematology analyzer, the Coulter Counter Model A, came out in 1956. The 1960s-1970s saw rapid development of multi-parameter analyzers. These could measure hemoglobin, WBC differential, and several red cell indices. The Technicon H-series (1970s-1980s) introduced cytochemical staining for differential counting. Modern analyzers (Sysmex XN, Beckman Coulter DxH, Siemens ADVIA) can measure more than 30 parameters from a single tube of blood. Turnaround times are under 1 minute. Automated morphology analysis and digital pathology are further transforming the CBC \cite{tietz2023textbook}.

\section*{Why This Test Is Ordered}
\begin{itemize}
  \item Routine health screening: the CBC is the most frequently ordered laboratory test in medicine, used for annual examinations and preoperative assessment
  \item Checking fatigue, weakness, pallor, or shortness of breath to screen for anemia
  \item Investigating suspected infection: fever, high white blood cell count, or signs of leukocytosis or leukopenia
  \item Assessing bruising, bleeding, or petechiae to evaluate platelet count and function
  \item Monitoring chemotherapy, radiation therapy, and immunosuppressive treatment, which suppress the bone marrow
  \item Monitoring chronic diseases: kidney disease, rheumatoid arthritis, inflammatory bowel disease, and connective tissue disorders
  \item Evaluating suspected hematologic malignancy: leukemia, lymphoma, or myelodysplastic syndromes
  \item Assessing bleeding risk before surgery or invasive procedures
  \item Monitoring response to treatment for anemia (iron, B12, folate) or infection
  \item Investigating lymphadenopathy, unexplained weight loss, or night sweats
  \item Screening for polycythemia or thrombocytosis in patients with suggestive symptoms or risk factors
\end{itemize}

\section*{Primary vs Secondary Test}
This is a primary screening test for a wide range of hematologic and systemic conditions.

\section*{How This Test Is Performed}
For most patients, a health care professional takes a blood sample from a vein in your arm using a small needle. After the needle goes in, a small amount of blood is collected into a test tube or vial. You may feel a little sting when the needle goes in or out. This usually takes less than five minutes. For certain tests, a urine sample, stool sample, or other specimen may be required \cite{clsi2023laboratory}.

\section*{What Preparation Is Needed}
Preparation requirements vary by test. Some tests require fasting for 8 to 12 hours. Others have no special preparation. Your healthcare provider or laboratory will tell you about any specific preparation needed for this test.

\section*{Contraindications and Precautions}
\begin{itemize}
  \item There are no absolute contraindications. Factors that can affect results: recent exercise, high altitude, dehydration (can raise hemoglobin/hematocrit), pregnancy (physiological hemodilution), recent blood transfusion, IV fluids. Certain medications may affect results (chemotherapy, antibiotics, corticosteroids). The sample should be processed within 4 hours of collection or refrigerated at 2-8 degrees C for up to 24 hours. EDTA (purple-top tube) is the standard anticoagulant for CBC.
\end{itemize}
\section*{Risks}
The risk is minimal. Slight pain or bruising at the venipuncture site. Rare: hematoma, infection, vasovagal syncope.

\section*{What Happens After the Test}
Results are usually available within 1 to 2 hours. Abnormal results may prompt further tests. These include: Peripheral blood smear review (if automated flags suggest abnormal cell morphology), Reticulocyte count (if anemia is present, to assess bone marrow response), Iron studies (if microcytic anemia is suspected), Vitamin B12 and folate levels (if macrocytic anemia is suspected), Flow cytometry (if leukemia or lymphoma is suspected), Bone marrow biopsy (if significant cytopenias or abnormal cells are present). Your healthcare provider will interpret results along with your symptoms, medical history, and physical examination \cite{tierney2023current}.

\section*{Understanding Results}

\subsection*{Below Normal}
\begin{itemize}
  \item Low white blood cell count (leukopenia): May point to viral infection, bone marrow suppression, autoimmune disease, medication effect, or aplastic anemia.
  \item Neutropenia (<1,500/uL) increases infection risk.
  \item Lymphopenia may point to HIV infection or immunosuppression.
  \item Low hemoglobin/hematocrit (anemia): Sorted by MCV into microcytic (iron deficiency, thalassemia), normocytic (chronic disease, acute blood loss), or macrocytic (B12/folate deficiency, liver disease, myelodysplastic syndrome) \cite{fischbach2023manual}.
  \item Low platelet count (thrombocytopenia): May point to immune thrombocytopenic purpura (ITP), thrombotic thrombocytopenic purpura (TTP), disseminated intravascular coagulation (DIC), bone marrow failure, or medication effect.
  \item Platelets <10,000/uL carry significant bleeding risk.
\end{itemize}

\subsection*{Above Normal}
\begin{itemize}
  \item High white blood cell count (leukocytosis): Often caused by infection (especially bacterial), inflammation, stress response, smoking, medication effects (corticosteroids), or hematologic malignancy (leukemia).
  \item A very high WBC (>50,000/uL) raises concern for leukemia.
  \item High hemoglobin/hematocrit: May point to polycythemia vera, dehydration, high altitude, chronic lung disease, or testosterone use.
  \item High platelet count (thrombocytosis): Reactive (infection, inflammation, iron deficiency) or primary (essential thrombocythemia, polycythemia vera).
\end{itemize}

\section*{Normal Results}
\begin{itemize}
  \item \textbf{Hemoglobin:} Males: 13.5--17.5 g/dL; Females: 12.0--16.0 g/dL
  \item \textbf{Hematocrit:} Males: 38.3--48.6\%; Females: 35.5--44.9\%
  \item \textbf{White blood cell count:} 4,500--11,000/$\mu$L
  \item \textbf{Platelet count:} 150,000--400,000/$\mu$L
  \item \textbf{Mean corpuscular volume (MCV):} 80--100 fL
  \item \textbf{Red cell distribution width (RDW):} 11.5--14.5\%
  \item \textbf{Differential:} Neutrophils 40--70\%, Lymphocytes 20--40\%, Monocytes 2--8\%, Eosinophils 1--4\%, Basophils 0--1\%
\end{itemize}

\section*{Follow-up Tests}
\begin{itemize}
  \item \textbf{Peripheral blood smear:} Microscopic exam of blood cells to check morphology when automated flags are present or abnormal cells are suspected
  \item \textbf{Reticulocyte count:} Measures young red blood cells to check the bone marrow response to anemia. A high reticulocyte count suggests hemolysis or recovery from blood loss; a low count suggests bone marrow failure
  \item \textbf{Iron studies:} Serum iron, ferritin, TIBC, transferrin saturation --- to check for iron deficiency anemia
  \item \textbf{Vitamin B12 and folate:} When macrocytosis (MCV $>$ 100 fL) is present
  \item \textbf{Flow cytometry:} When leukemia or lymphoma is suspected based on abnormal cells or cytopenias
  \item \textbf{Bone marrow biopsy:} When significant cytopenias, unexplained abnormal cells, or suspected hematologic malignancy are present
\end{itemize}


\section*{References}
\bibliographystyle{unsrtnat}
\bibliography{references}

\clearpage
\section*{Regulatory Status and Authorization Date}
This chapter describes a test category, not a specific commercial product. FDA, CDSCO, EU IVDR, and MHRA approval or registration dates therefore cannot be assigned without the assay manufacturer, product name, instrument, and intended use. For laboratory-developed tests, an individual agency approval date may not exist. See the Preface for the interpretation of regulatory status.

\clearpage
